\documentclass[a4paper]{article}
\usepackage{graphicx}
\usepackage{amsmath,amssymb}
\usepackage{minted}
\usepackage{multirow}
\usepackage{float}
\usepackage{algorithm}
\usepackage{algpseudocode}
\usepackage{todonotes}
\usepackage{forest}
\usepackage{xstring}
\usepackage{xcolor}
\usepackage[most]{tcolorbox}
\usepackage{xparse}
\usepackage{natbib}
\usetikzlibrary{graphs,graphdrawing}
\usegdlibrary{layered}
% ND Mind Mirror default input files.
\input{/home/donkarlo/Dropbox/repo/nd_language_ai_project/src/nd_language_ai/formal/computer/programming/tex/latex/code_clips/external_graphviz.tex}
\input{/home/donkarlo/Dropbox/repo/nd_language_ai_project/src/nd_language_ai/formal/computer/programming/tex/latex/part/control_sequence/kind/semtag/semtag.tex}
\input{/home/donkarlo/Dropbox/repo/nd_language_ai_project/src/nd_language_ai/formal/computer/programming/tex/latex/code_clips/seehere.tex}

\title{STL Trend and Seasonal Decomposition}
\date{}
\author{Mohammad Rahmani}
\begin{document}
\maketitle

\section{Definition}
STL means \emph{Seasonal and Trend decomposition using LOESS}. It decomposes regularly observed data into a time-varying trend, recurring seasonal effects, and a remainder.

\section{Additive Model}
\[
 y_t=T_t+S_t+R_t.
\]
Here $T_t$ is a smooth trend, $S_t$ is the seasonal signal, and $R_t$ captures what the first two components do not explain. Additive STL is less natural for series whose seasonal amplitude is proportional to their level; a transformation such as logarithms may be useful for strictly positive observations.

\section{Methodology}
Given observations with a consistent frequency and a seasonal period $p$:
\begin{enumerate}
 \item Smooth groups of observations at corresponding seasonal positions.
 \item Construct an estimate of the seasonal component.
 \item Remove seasonality and locally smooth the deseasonalized data to estimate trend.
 \item Iterate estimates of the seasonal and trend components.
 \item Optionally downweight outliers with robust iterations.
 \item Compute residuals $R_t=y_t-T_t-S_t$.
\end{enumerate}
LOESS smoothing is a central ingredient of the decomposition.

\section{E-Commerce Example}
Daily e-commerce orders commonly differ between weekdays and weekends. With seasonal period $p=7$, STL can help isolate an underlying growth pattern from these weekly changes.

\section{Python Illustration}
\begin{verbatim}
from statsmodels.tsa.seasonal import STL

# daily_sales: a regular daily Series, with missing
# calendar dates and missing values handled explicitly.
fit = STL(daily_sales, period=7, robust=True).fit()
trend = fit.trend
seasonality = fit.seasonal
remainder = fit.resid
\end{verbatim}

\section{Advantages}
\begin{itemize}
 \item Separates long-run changes from a recurring seasonal pattern.
 \item Supports nonlinear smooth trends.
 \item Offers robust fitting to limit certain outlier effects.
\end{itemize}

\section{Limitations}
\begin{itemize}
 \item Requires appropriate frequency and seasonal-period selection.
 \item Standard STL uses one seasonal period; use a multiple-seasonality approach for weekly plus annual patterns.
 \item Trend estimates depend on smoothing settings and may be unstable near endpoints.
 \item Retrospective two-sided smoothing may use future observations and leak information in predictive modeling.
\end{itemize}

\section{Practical Guidance}
Inspect all three components. Account for promotions, holidays, stockouts, and known structural changes. During chronological forecasting evaluation, fit preprocessing only to the available training history.

\bibliographystyle{plainnat}
\bibliography{/home/donkarlo/Dropbox/repo/nd_shared_research/bibliography/bibliography.bib}
\end{document}
